\chapter[Referencial teórico]{Referencial Teórico}

Este capítulo apresenta a fundamentação teórica que sustenta esse trabalho, trazendo a tona os principais 
conceitos, teorias e estudos relacionados ao tema abordado.

\section{Engenharia de Software Experimental}

Diferente de áreas científicas mais precisas, como a física, onde fenômenos podem ser modelados 
com rigor matemático, a engenharia de software é uma disciplina profundamente influenciada pela 
complexidade do comportamento humano e pelas interações dinâmicas que ocorrem durante o processo 
de desenvolvimento, com decisões e práticas que variam conforme o contexto, a equipe envolvida e 
as ferramentas utilizadas. Por isso, a experimentação se torna uma abordagem essencial para 
garantir que escolhas metodológicas, técnicas e ferramentas sejam fundamentadas em dados empíricos. 
A Engenharia de Software Experimental proporciona os meios para testar e validar abordagens, 
possibilitando uma avaliação científica das práticas adotadas no desenvolvimento de sistemas \cite{Wohlin:2012:ESE:2349018}.

Entre os métodos de pesquisa no contexto da Engenharia de Software apresentados por \citeonline{Basili1993}, 
destaca-se o método empírico. Esse método baseia-se na observação direta e na coleta de dados reais, buscando validar teorias 
e modelos por meio da análise de dados obtidos de experimentos práticos. A pesquisa empírica é amplamente utilizada na ciencias sociais onde se tem foco no comportamento humano,
em que não podemos estabelecer leis da natureza. Portanto tornando-se essencial em um campo como a engenharia de software, que envolve a interação de múltiplos fatores humanos, 
tecnológicos e organizacionais. \cite{Wohlin:2012:ESE:2349018}.

O Experimento é uma das principais estratégias de pesquisa empírica. Ele envolve a manipulação controlada de fatores ou variáveis.
Diferentes tratamentos são aplicados a ou por diferentes sujeitos, medindo os efeitos nas variáveis de resultado. É especialmente
utilizado em estudos onde se deseja comparar o desempenho de diferentes técnicas, ferramentas ou metodologias de desenvolvimento de software \cite{Wohlin:2012:ESE:2349018}.


\subsection{Experimento}

O experimento é uma abordagem para investigar hipóteses por meio da manipulação controlada de variáveis.
Ele é projetado para estabelecer relações de causa e efeito, permitindo que os pesquisadores isolem o impacto de uma ou mais variáveis 
independentes sobre uma ou mais variáveis dependentes. Segundo \citeonline{Wohlin:2012:ESE:2349018}, um experimento envolve as seguintes etapas principais:

\subsubsection{Definição do Escopo}
Nesta primeira etapa, a fundação do experimento é determinada, e o experimento é delimitado em termos de problema, objetivo e metas.
O objetivo é garantir que a intenção do experimento possa ser cumprida.
A meta do experimento é formalizada usando uma estrutura que define: Objeto de estudo, Propósito, Foco de qualidade, Perspectiva e Contexto.

\subsubsection{Planejamento (Planning)}
Esta fase prepara o "como" o experimento será conduzido.
O planejamento estabelece a fundação para o experimento.
As principais sub-etapas incluem:
\begin{itemize}

	\item Seleção do Contexto: Determinação do ambiente onde o experimento será executado . 
	\item Formulação de Hipóteses: Onde a hipótese é formalizada, incluindo a hipótese nula. 
	\item Seleção de Variáveis: Definição das variáveis dependentes (o que será medido) e independentes (o que será manipulado).
	\item Seleção dos Sujeitos: Escolha da amostra da população de interesse.
	\item Escolha do Tipo de Design: Definição do design experimental apropriado.
	\item Instrumentação: Preparação dos instrumentos necessários, como objetos de experimento, diretrizes e instrumentos de medição.
	\item Avaliação da Validade: Análise dos riscos (ameaças) à validade dos resultados em quatro categorias principais: conclusão, interna, construto e externa.

\end{itemize}

\subsubsection{Operação (Operation)}
Fase de execução do experimento, onde os tratamentos são aplicados aos sujeitos e os dados são coletados.
É crucial que o experimento seja conduzido de acordo com o plano para que os resultados sejam válidos.
As principais sub-etapas incluem:

\begin{itemize}
    \item Preparação: Seleção e informação dos participantes (garantindo o consentimento e a ética) e preparação de materiais, como formulários e ferramentas.
    \item Execução: Os sujeitos realizam suas tarefas conforme os diferentes tratamentos, e os dados são coletados (manualmente, com ferramentas ou automaticamente).
    \item Validação de Dados: Verificação dos dados coletados para garantir que sejam razoáveis e corretos, e identificação de outliers ou 
	de dados inválidos que precisam ser removidos.

\end{itemize}

\subsubsection{Análise e Interpretação (Analysis and Interpretation)}
Nesta fase, os dados coletados são processados para que as conclusões possam ser tiradas.
As etapas de interpretação quantitativa são:
\begin{itemize}
    \item Estatísticas Descritivas: Caracterização dos dados e visualização gráfica.
    \item Redução do Conjunto de Dados: Exclusão de pontos de dados anormais ou falsos (como outliers), para obter um conjunto de dados válidos para a análise.
    \item Teste de Hipóteses: Avaliação estatística das hipóteses do experimento, utilizando testes paramétricos ou não-paramétricos para ver se a 
	hipótese nula pode ser rejeitada com um nível de significância definido.
    \item Desenho de Conclusões: Determinação se a hipótese pode ser 
	rejeitada, formando a base para decisões e conclusões sobre como usar os resultados.
\end{itemize}

\subsubsection{Apresentação e Empacotamento (Presentation and Package)}
Esta atividade final envolve a documentação e comunicação das descobertas.
O foco principal é a documentação dos resultados, seja em um artigo de pesquisa, um relatório para tomadores de decisão, ou 
como parte da base de experiência de uma empresa. É essencial para o aprendizado e para facilitar a replicação do experimento 
por outros pesquisadores.


\section{Modelos de Qualidade}

O Desenvolvimento (Miolo ou Corpo do Trabalho) é subdividido em seções de 
acordo com o planejamento do autor. As seções primárias são aquelas que 
resultam da primeira divisão do texto do documento, geralmente 
correspondendo a divisão em capítulos. Seções secundárias, terciárias, 
etc., são aquelas que resultam da divisão do texto de uma seção primária, 
secundária, terciária, etc., respectivamente.

As seções primárias são numeradas consecutivamente, seguindo a série 
natural de números inteiros, a partir de 1, pela ordem de sua sucessão no 
documento.

O Desenvolvimento é a seção mais importante do trabalho, por isso exigi-se 
organização, objetividade e clareza. É conveniente dividi-lo em pelo menos 
três partes:

\begin{itemize}

	\item Referencial teórico, que corresponde a uma análise dos trabalhos 
	relevantes, encontrados na pesquisa bibliográfica sobre o assunto. 
	\item Metodologia, que é a descrição de todos os passos metodológicos 
	utilizados no trabalho. Sugere-se que se enfatize especialmente em (1) 
	População ou Sujeitos da pesquisa, (2) Materiais e equipamentos 
	utilizados e (3) Procedimentos de coleta de dados.
	\item Resultados, Discussão dos resultados e Conclusões, que é onde se 
	apresenta os dados encontrados a análise feita pelo autor à luz do 
	Referencial teórico e as Conclusões.

\end{itemize}

\section{Metodologias Ágeis}

O uso de programas de edição eletrônica de textos é de livre escolha do autor. 

\section{Code Smells}

O uso de programas de edição eletrônica de textos é de livre escolha do autor. 

\section{Inteligência Artificial}

O uso de programas de edição eletrônica de textos é de livre escolha do autor. 

